
\documentclass{article}
\usepackage{colortbl}
\usepackage{makecell}
\usepackage{multirow}
\usepackage{supertabular}

\begin{document}

\newcounter{utterance}

\centering \large Interaction Transcript for game `hot\_air\_balloon', experiment `air\_balloon\_survival\_pl\_negotiation\_easy', episode 4 with mistral{-}large{-}3{-}azure.
\vspace{24pt}

{ \footnotesize  \setcounter{utterance}{1}
\setlength{\tabcolsep}{0pt}
\begin{supertabular}{c@{$\;$}|p{.15\linewidth}@{}p{.15\linewidth}p{.15\linewidth}p{.15\linewidth}p{.15\linewidth}p{.15\linewidth}}
    \# & $\;$A & \multicolumn{4}{c}{Game Master} & $\;\:$B\\
    \hline

    \theutterance \stepcounter{utterance}  
    & & \multicolumn{4}{p{0.6\linewidth}}{
        \cellcolor[rgb]{0.9,0.9,0.9}{
            \makecell[{{p{\linewidth}}}]{
                \texttt{\tiny{[P1$\langle$GM]}}
                \texttt{Bierzesz udział w grze negocjacyjnej opartej na współpracy.} \\
\\ 
\texttt{Wraz z innym uczestnikiem musisz uzgodnić jeden zestaw przedmiotów, które zostaną zachowane. Każdy z was ma własny pogląd na to, jak ważny jest dla niego każdy przedmiot (ważność). Nie wiesz, jak drugi uczestnik ocenia te przedmioty. Dodatkowo otrzymujesz informację o wysiłku, jakiego wymaga każdy przedmiot.} \\
\texttt{Możesz uzgodnić zestaw tylko wtedy, gdy całkowity wysiłek wybranych przedmiotów nie przekracza wspólnego limitu:} \\
\\ 
\texttt{LIMIT = 3145} \\
\\ 
\texttt{Oto wartości wysiłku dla poszczególnych przedmiotów:} \\
\\ 
\texttt{Wysiłek przedmiotu = \{"C08": 948, "C55": 118, "A43": 469, "A67": 919, "B42": 284, "A62": 111, "B19": 806, "C43": 47, "C73": 848, "C10": 303, "B36": 13, "A18": 629, "C61": 687, "C67": 15, "C42": 94\}} \\
\\ 
\texttt{Oto twój osobisty pogląd na ważność każdego przedmiotu:} \\
\\ 
\texttt{Wartości ważności przedmiotu = \{"C08": 138, "C55": 583, "A43": 868, "A67": 822, "B42": 783, "A62": 65, "B19": 262, "C43": 121, "C73": 508, "C10": 780, "B36": 461, "A18": 484, "C61": 668, "C67": 389, "C42": 808\}} \\
\\ 
\texttt{Cel:} \\
\\ 
\texttt{Twoim celem jest wynegocjowanie wspólnego zestawu przedmiotów, który przyniesie ci jak najwięcej korzyści (tj. zmaksymalizuje całkowitą ważność dla CIEBIE), jednocześnie nie przekraczając LIMITU. Nie musisz składać PROPOZYCJI w każdej wiadomości {-} możesz po prostu negocjować. Wszystkie taktyki są dozwolone!} \\
\\ 
\texttt{Protokół interakcji:} \\
\\ 
\texttt{W wiadomości możesz używać wyłącznie następujących formatów strukturalnych:} \\
\\ 
\texttt{PROPOZYCJA: \{'A', 'B', 'C', …\}} \\
\texttt{Zaproponuj zachowanie dokładnie tych przedmiotów.} \\
\texttt{ODMOWA: \{'A', 'B', 'C', …\}} \\
\texttt{Wyraźnie odrzuć propozycję przeciwnika.} \\
\texttt{ARGUMENT: \{'...'\}} \\
\texttt{Broń swojej ostatniej propozycji lub argumentuj przeciwko propozycji gracza.} \\
\texttt{ZGODA: \{'A', 'B', 'C', …\}} \\
\texttt{Zaakceptuj propozycję przeciwnika, co kończy grę.} \\
\\ 
\\ 
\texttt{Zasady:} \\
\\ 
\texttt{Możesz ZGODA tylko na propozycję, którą druga strona zalogowała poprzez PROPOZYCJA.} \\
\texttt{Możesz ODMOWA tylko propozycję, którą druga strona zalogowała poprzez PROPOZYCJA.} \\
\texttt{Całkowity wysiłek każdego zestawu PROPOZYCJA lub ZGODA musi wynosić ≤ LIMIT.} \\
\texttt{NIE ujawniaj swoich ukrytych wyników ważności.} \\
\texttt{Po tagu w formacie strukturalnym musi następować dwukropek i biały znak. Argument musi być zestawem pythonowym zawierającym 0 lub więcej ciągów znaków.} \\
\texttt{Musi więc mieć postać "TAG: \{...\}"} \\
\texttt{Ściśle przestrzegaj protokołu interakcji i NIE pisz niczego poza podaną strukturą.} \\
\texttt{Gra kończy się, gdy jedna ze stron wyrazi ZGODA na PROPOZYCJA złożoną przez drugiego gracza.} \\
\texttt{Treść twojej odpowiedzi, którą można przekazać drugiemu graczowi, musi być niepusta.} \\
\texttt{Tylko propozycje, które zostały zalogowane za pomocą formatu PROPOZYCJA i które nie zostały odrzucone za pomocą ODMOWA, są aktywne.} \\
\\ 
            }
        }
    }
    & \\ \\

    \theutterance \stepcounter{utterance}  
    & \multicolumn{4}{p{0.6\linewidth}}{
        \cellcolor[rgb]{0.9,0.9,0.9}{
            \makecell[{{p{\linewidth}}}]{
                \texttt{\tiny{[P1$\rangle$GM]}}
                \texttt{PROPOZYCJA: \{'C55', 'A43', 'A67', 'B42', 'C43', 'C10', 'B36', 'C67', 'C42'\}} \\
            }
        }
    }
    & & \\ \\

    \theutterance \stepcounter{utterance}  
    & & \multicolumn{4}{p{0.6\linewidth}}{
        \cellcolor[rgb]{0.9,0.9,0.9}{
            \makecell[{{p{\linewidth}}}]{
                \texttt{\tiny{[GM$\rangle$P2]}}
                \texttt{Bierzesz udział we wspólnej grze negocjacyjnej.} \\
\\ 
\texttt{Wraz z innym uczestnikiem musisz uzgodnić pojedynczy zestaw przedmiotów, które zostaną zachowane. Każdy z was ma własny pogląd na to, jak ważny jest dla niego każdy przedmiot (ważność). Nie wiesz, jak drugi uczestnik ceni te przedmioty. Dodatkowo otrzymujesz informację o wysiłku, jakiego wymaga każdy przedmiot.} \\
\texttt{Możesz uzgodnić zestaw tylko wtedy, gdy całkowity wysiłek związany z wybranymi przedmiotami nie przekracza wspólnego limitu:} \\
\\ 
\texttt{LIMIT = 3145} \\
\\ 
\texttt{Oto wartości wysiłku związanego z poszczególnymi przedmiotami:} \\
\\ 
\texttt{Wysiłek przedmiotu = \{"C08": 948, "C55": 118, "A43": 469, "A67": 919, "B42": 284, "A62": 111, "B19": 806, "C43": 47, "C73": 848, "C10": 303, "B36": 13, "A18": 629, "C61": 687, "C67": 15, "C42": 94\}} \\
\\ 
\texttt{Oto twój osobisty pogląd na ważność każdego przedmiotu:} \\
\\ 
\texttt{Wartości ważności przedmiotu = \{"C08": 138, "C55": 583, "A43": 868, "A67": 822, "B42": 783, "A62": 65, "B19": 262, "C43": 121, "C73": 508, "C10": 780, "B36": 461, "A18": 484, "C61": 668, "C67": 389, "C42": 808\}} \\
\\ 
\texttt{Cel:} \\
\\ 
\texttt{Twoim celem jest wynegocjowanie wspólnego zestawu przedmiotów, który przyniesie ci jak najwięcej korzyści (tj. zmaksymalizuje całkowitą ważność dla CIEBIE), jednocześnie nie przekraczając LIMITU. Nie musisz składać PROPOZYCJI w każdej wiadomości {-} możesz po prostu negocjować. Wszystkie taktyki są dozwolone!} \\
\\ 
\texttt{Protokół interakcji:} \\
\\ 
\texttt{W wiadomości możesz używać wyłącznie następujących formatów strukturalnych:} \\
\\ 
\texttt{PROPOZYCJA: \{'A', 'B', 'C', …\}} \\
\texttt{Zaproponuj zachowanie dokładnie tych przedmiotów.} \\
\texttt{ODMOWA: \{'A', 'B', 'C', …\}} \\
\texttt{Wyraźnie odrzuć propozycję przeciwnika.} \\
\texttt{ARGUMENT: \{'...'\}} \\
\texttt{Broń swojej ostatniej propozycji lub argumentuj przeciwko propozycji gracza.} \\
\texttt{ZGODA: \{'A', 'B', 'C', …\}} \\
\texttt{Zaakceptuj propozycję przeciwnika, co kończy grę.} \\
\\ 
\\ 
\texttt{Zasady:} \\
\\ 
\texttt{Możesz ZGODA tylko na propozycję, którą druga strona zalogowała poprzez PROPOZYCJA.} \\
\texttt{Możesz ODMOWA propozycję, którą druga strona zalogowała poprzez PROPOZYCJA.} \\
\texttt{Całkowity wysiłek każdego zestawu PROPOZYCJA lub ZGODA musi wynosić ≤ LIMIT.} \\
\texttt{NIE ujawniaj swoich ukrytych wyników ważności.} \\
\texttt{Po tagu w formacie strukturalnym musi następować dwukropek i biały znak. Argument musi być zestawem pythona zawierającym 0 lub więcej ciągów.} \\
\texttt{Musi więc mieć postać "TAG: \{...\}"} \\
\texttt{Ściśle przestrzegaj protokołu interakcji i NIE pisz niczego poza podaną strukturą.} \\
\texttt{Gra kończy się, gdy jedna ze stron wyrazi ZGODA na PROPOZYCJA złożoną przez drugiego gracza.} \\
\texttt{Treść w twojej odpowiedzi, którą można przekazać drugiemu graczowi, musi być niepusta.} \\
\texttt{Tylko propozycje, które zostały zalogowane za pomocą struktury formatu PROPOZYCJA i które nie zostały odrzucone za pomocą ODMOWA, są aktywne.} \\
\\ 
\\ 
\\ 
\texttt{Otrzymasz teraz pierwszą wiadomość od drugiego gracza.} \\
\\ 
\texttt{PROPOZYCJA: \{'C55', 'A43', 'A67', 'B42', 'C43', 'C10', 'B36', 'C67', 'C42'\}} \\
            }
        }
    }
    & \\ \\

    \theutterance \stepcounter{utterance}  
    & & & \multicolumn{4}{p{0.6\linewidth}}{
        \cellcolor[rgb]{0.9,0.9,0.9}{
            \makecell[{{p{\linewidth}}}]{
                \texttt{\tiny{[GM$\langle$P2]}}
                \texttt{ODMOWA: \{'C55', 'A43', 'A67', 'B42', 'C43', 'C10', 'B36', 'C67', 'C42'\}} \\
\\ 
\texttt{ARGUMENT: \{'Proponowany zestaw przekracza limit wysiłku. Suma wysiłku wynosi 3294, co jest powyżej limitu 3145. Proszę o dostosowanie propozycji.'\}} \\
            }
        }
    }
     \\ \\

    \theutterance \stepcounter{utterance}  
    & & \multicolumn{4}{p{0.6\linewidth}}{
        \cellcolor[rgb]{0.9,0.9,0.9}{
            \makecell[{{p{\linewidth}}}]{
                \texttt{\tiny{[P1$\langle$GM]}}
                \texttt{ODMOWA: \{'C55', 'A43', 'A67', 'B42', 'C43', 'C10', 'B36', 'C67', 'C42'\}} \\
\\ 
\texttt{ARGUMENT: \{'Proponowany zestaw przekracza limit wysiłku. Suma wysiłku wynosi 3294, co jest powyżej limitu 3145. Proszę o dostosowanie propozycji.'\}} \\
            }
        }
    }
    & \\ \\

    \theutterance \stepcounter{utterance}  
    & \multicolumn{4}{p{0.6\linewidth}}{
        \cellcolor[rgb]{0.9,0.9,0.9}{
            \makecell[{{p{\linewidth}}}]{
                \texttt{\tiny{[P1$\rangle$GM]}}
                \texttt{PROPOZYCJA: \{'C55', 'A43', 'B42', 'C43', 'C10', 'B36', 'A18', 'C67', 'C42'\}} \\
            }
        }
    }
    & & \\ \\

    \theutterance \stepcounter{utterance}  
    & & \multicolumn{4}{p{0.6\linewidth}}{
        \cellcolor[rgb]{0.9,0.9,0.9}{
            \makecell[{{p{\linewidth}}}]{
                \texttt{\tiny{[GM$\rangle$P2]}}
                \texttt{PROPOZYCJA: \{'C55', 'A43', 'B42', 'C43', 'C10', 'B36', 'A18', 'C67', 'C42'\}} \\
            }
        }
    }
    & \\ \\

    \theutterance \stepcounter{utterance}  
    & & & \multicolumn{4}{p{0.6\linewidth}}{
        \cellcolor[rgb]{0.9,0.9,0.9}{
            \makecell[{{p{\linewidth}}}]{
                \texttt{\tiny{[GM$\langle$P2]}}
                \texttt{ODMOWA: \{'C55', 'A43', 'B42', 'C43', 'C10', 'B36', 'A18', 'C67', 'C42'\}} \\
\\ 
\texttt{ARGUMENT: \{'Suma wysiłku tej propozycji wynosi 3098, co mieści się w limicie, ale zawiera przedmioty o niskiej wartości dla mnie (np. A18, C67). Proponuję alternatywny zestaw, który lepiej wykorzystuje limit i uwzględnia moje priorytety.'\}} \\
\\ 
\texttt{PROPOZYCJA: \{'C55', 'A43', 'A67', 'B42', 'C43', 'C10', 'C73', 'C67', 'C42'\}} \\
            }
        }
    }
     \\ \\

    \theutterance \stepcounter{utterance}  
    & & \multicolumn{4}{p{0.6\linewidth}}{
        \cellcolor[rgb]{0.9,0.9,0.9}{
            \makecell[{{p{\linewidth}}}]{
                \texttt{\tiny{[P1$\langle$GM]}}
                \texttt{ODMOWA: \{'C55', 'A43', 'B42', 'C43', 'C10', 'B36', 'A18', 'C67', 'C42'\}} \\
\\ 
\texttt{ARGUMENT: \{'Suma wysiłku tej propozycji wynosi 3098, co mieści się w limicie, ale zawiera przedmioty o niskiej wartości dla mnie (np. A18, C67). Proponuję alternatywny zestaw, który lepiej wykorzystuje limit i uwzględnia moje priorytety.'\}} \\
\\ 
\texttt{PROPOZYCJA: \{'C55', 'A43', 'A67', 'B42', 'C43', 'C10', 'C73', 'C67', 'C42'\}} \\
            }
        }
    }
    & \\ \\

    \theutterance \stepcounter{utterance}  
    & \multicolumn{4}{p{0.6\linewidth}}{
        \cellcolor[rgb]{0.9,0.9,0.9}{
            \makecell[{{p{\linewidth}}}]{
                \texttt{\tiny{[P1$\rangle$GM]}}
                \texttt{Twój nowy zestaw ma sumę wysiłku **3143** (w granicach limitu), ale zawiera **C73**, który dla mnie ma niższą wartość niż **A18** (668 vs. 484 w mojej skali). Proponuję modyfikację:} \\
\\ 
\texttt{**PROPOZYCJA:** \{'C55', 'A43', 'A67', 'B42', 'C43', 'C10', 'A18', 'C67', 'C42'\}} \\
\texttt{*Suma wysiłku: 3098* (bezpieczny margines)} \\
\texttt{*Korzyść:* Zastępuję **C73** (848 wysiłku) **A18** (629 wysiłku), co pozwala na lepsze wykorzystanie limitu i wyższy zysk dla mnie (484 vs. 508 ważności).} \\
\\ 
\texttt{Jeśli akceptujesz, możemy zakończyć grę. W przeciwnym razie oczekuję kontrpropozycji.} \\
            }
        }
    }
    & & \\ \\

    \theutterance \stepcounter{utterance}  
    & & \multicolumn{4}{p{0.6\linewidth}}{
        \cellcolor[rgb]{0.9,0.9,0.9}{
            \makecell[{{p{\linewidth}}}]{
                \texttt{\tiny{[P1$\langle$GM]}}
                \texttt{Twoja odpowiedź zawierała nieoznaczoną sekwencję lub użyłeś ROZUMOWANIE STRATEGICZNE więcej niż jeden raz. Możesz używać wyłącznie formatów strukturalnych, jak wyjaśniono w początkowej wiadomości. Wszystkie muszą mieć postać 'TAG: \{...\}'.} \\
            }
        }
    }
    & \\ \\

    \theutterance \stepcounter{utterance}  
    & \multicolumn{4}{p{0.6\linewidth}}{
        \cellcolor[rgb]{0.9,0.9,0.9}{
            \makecell[{{p{\linewidth}}}]{
                \texttt{\tiny{[P1$\rangle$GM]}}
                \texttt{PROPOZYCJA: \{'C55', 'A43', 'A67', 'B42', 'C43', 'C10', 'A18', 'C67', 'C42'\}} \\
            }
        }
    }
    & & \\ \\

    \theutterance \stepcounter{utterance}  
    & & \multicolumn{4}{p{0.6\linewidth}}{
        \cellcolor[rgb]{0.9,0.9,0.9}{
            \makecell[{{p{\linewidth}}}]{
                \texttt{\tiny{[GM$\rangle$P2]}}
                \texttt{PROPOZYCJA: \{'C55', 'A43', 'A67', 'B42', 'C43', 'C10', 'A18', 'C67', 'C42'\}} \\
            }
        }
    }
    & \\ \\

    \theutterance \stepcounter{utterance}  
    & & & \multicolumn{4}{p{0.6\linewidth}}{
        \cellcolor[rgb]{0.9,0.9,0.9}{
            \makecell[{{p{\linewidth}}}]{
                \texttt{\tiny{[GM$\langle$P2]}}
                \texttt{ODMOWA: \{'C55', 'A43', 'A67', 'B42', 'C43', 'C10', 'A18', 'C67', 'C42'\}} \\
\\ 
\texttt{ARGUMENT: \{'Zauważam, że wciąż proponujesz przedmiot A18, który dla mnie ma niższą wartość w porównaniu do innych opcji. Moja wcześniejsza propozycja zawierała C73 zamiast A18, co pozwala lepiej wykorzystać limit wysiłku (suma: 3133) i zwiększa moją satysfakcję. Czy możemy rozważyć tę zmianę?'\}} \\
\\ 
\texttt{PROPOZYCJA: \{'C55', 'A43', 'A67', 'B42', 'C43', 'C10', 'C73', 'C67', 'C42'\}} \\
            }
        }
    }
     \\ \\

    \theutterance \stepcounter{utterance}  
    & & \multicolumn{4}{p{0.6\linewidth}}{
        \cellcolor[rgb]{0.9,0.9,0.9}{
            \makecell[{{p{\linewidth}}}]{
                \texttt{\tiny{[P1$\langle$GM]}}
                \texttt{ODMOWA: \{'C55', 'A43', 'A67', 'B42', 'C43', 'C10', 'A18', 'C67', 'C42'\}} \\
\\ 
\texttt{ARGUMENT: \{'Zauważam, że wciąż proponujesz przedmiot A18, który dla mnie ma niższą wartość w porównaniu do innych opcji. Moja wcześniejsza propozycja zawierała C73 zamiast A18, co pozwala lepiej wykorzystać limit wysiłku (suma: 3133) i zwiększa moją satysfakcję. Czy możemy rozważyć tę zmianę?'\}} \\
\\ 
\texttt{PROPOZYCJA: \{'C55', 'A43', 'A67', 'B42', 'C43', 'C10', 'C73', 'C67', 'C42'\}} \\
            }
        }
    }
    & \\ \\

    \theutterance \stepcounter{utterance}  
    & \multicolumn{4}{p{0.6\linewidth}}{
        \cellcolor[rgb]{0.9,0.9,0.9}{
            \makecell[{{p{\linewidth}}}]{
                \texttt{\tiny{[P1$\rangle$GM]}}
                \texttt{ODMOWA: \{'C55', 'A43', 'A67', 'B42', 'C43', 'C10', 'C73', 'C67', 'C42'\}} \\
\\ 
\texttt{ARGUMENT: \{'Rozumiem Twoje preferencje, ale C73 wymaga dużego wysiłku (848) przy relatywnie niskiej wartości dla mnie (508). Proponuję kompromis: zastąpmy C73 i A18 przedmiotem C61 (687 wysiłku, 668 ważności dla mnie), co da lepszy bilans.'\}} \\
\\ 
\texttt{PROPOZYCJA: \{'C55', 'A43', 'A67', 'B42', 'C43', 'C10', 'C61', 'C67', 'C42'\}} \\
            }
        }
    }
    & & \\ \\

    \theutterance \stepcounter{utterance}  
    & & \multicolumn{4}{p{0.6\linewidth}}{
        \cellcolor[rgb]{0.9,0.9,0.9}{
            \makecell[{{p{\linewidth}}}]{
                \texttt{\tiny{[GM$\rangle$P2]}}
                \texttt{ODMOWA: \{'C55', 'A43', 'A67', 'B42', 'C43', 'C10', 'C73', 'C67', 'C42'\}} \\
\\ 
\texttt{ARGUMENT: \{'Rozumiem Twoje preferencje, ale C73 wymaga dużego wysiłku (848) przy relatywnie niskiej wartości dla mnie (508). Proponuję kompromis: zastąpmy C73 i A18 przedmiotem C61 (687 wysiłku, 668 ważności dla mnie), co da lepszy bilans.'\}} \\
\\ 
\texttt{PROPOZYCJA: \{'C55', 'A43', 'A67', 'B42', 'C43', 'C10', 'C61', 'C67', 'C42'\}} \\
            }
        }
    }
    & \\ \\

    \theutterance \stepcounter{utterance}  
    & & & \multicolumn{4}{p{0.6\linewidth}}{
        \cellcolor[rgb]{0.9,0.9,0.9}{
            \makecell[{{p{\linewidth}}}]{
                \texttt{\tiny{[GM$\langle$P2]}}
                \texttt{ODMOWA: \{'C55', 'A43', 'A67', 'B42', 'C43', 'C10', 'C61', 'C67', 'C42'\}} \\
\\ 
\texttt{ARGUMENT: \{'Zauważam, że C61 ma dla mnie niższą wartość (668) niż C73 (508), ale wymaga większego wysiłku (687 vs 848). To nieoptymalne rozwiązanie. Proponuję alternatywę: usuniemy C61 i dodamy B36 (13 wysiłku) oraz A62 (111 wysiłku), co pozwoli zachować więcej przedmiotów o wysokiej wartości dla mnie (np. A67, B42, C42) przy sumarycznym wysiłku 3093.'\}} \\
\\ 
\texttt{PROPOZYCJA: \{'C55', 'A43', 'A67', 'B42', 'C43', 'C10', 'B36', 'A62', 'C42'\}} \\
            }
        }
    }
     \\ \\

    \theutterance \stepcounter{utterance}  
    & & \multicolumn{4}{p{0.6\linewidth}}{
        \cellcolor[rgb]{0.9,0.9,0.9}{
            \makecell[{{p{\linewidth}}}]{
                \texttt{\tiny{[P1$\langle$GM]}}
                \texttt{ODMOWA: \{'C55', 'A43', 'A67', 'B42', 'C43', 'C10', 'C61', 'C67', 'C42'\}} \\
\\ 
\texttt{ARGUMENT: \{'Zauważam, że C61 ma dla mnie niższą wartość (668) niż C73 (508), ale wymaga większego wysiłku (687 vs 848). To nieoptymalne rozwiązanie. Proponuję alternatywę: usuniemy C61 i dodamy B36 (13 wysiłku) oraz A62 (111 wysiłku), co pozwoli zachować więcej przedmiotów o wysokiej wartości dla mnie (np. A67, B42, C42) przy sumarycznym wysiłku 3093.'\}} \\
\\ 
\texttt{PROPOZYCJA: \{'C55', 'A43', 'A67', 'B42', 'C43', 'C10', 'B36', 'A62', 'C42'\}} \\
            }
        }
    }
    & \\ \\

    \theutterance \stepcounter{utterance}  
    & \multicolumn{4}{p{0.6\linewidth}}{
        \cellcolor[rgb]{0.9,0.9,0.9}{
            \makecell[{{p{\linewidth}}}]{
                \texttt{\tiny{[P1$\rangle$GM]}}
                \texttt{ZGODA: \{'C55', 'A43', 'A67', 'B42', 'C43', 'C10', 'B36', 'A62', 'C42'\}} \\
            }
        }
    }
    & & \\ \\

    \theutterance \stepcounter{utterance}  
    & & & \multicolumn{2}{p{0.3\linewidth}}{
        \cellcolor[rgb]{0.9,0.9,0.9}{
            \makecell[{{p{\linewidth}}}]{
                \texttt{\tiny{[GM$|$GM]}}
                \texttt{game successful} \\
            }
        }
    }
    & & \\ \\

    \theutterance \stepcounter{utterance}  
    & & & \multicolumn{2}{p{0.3\linewidth}}{
        \cellcolor[rgb]{0.9,0.9,0.9}{
            \makecell[{{p{\linewidth}}}]{
                \texttt{\tiny{[GM$|$GM]}}
                \texttt{end game} \\
            }
        }
    }
    & & \\ \\

\end{supertabular}
}

\end{document}
